Small lesions in brain magnetic resonance imaging (MRI) are sparse,
spatially localized, and clinically important targets embedded within a
large volume of normal-appearing tissue. In multiple sclerosis (MS), lesions
occupy little of the image and vary widely in size, yet missing them can
underestimate disease activity and lesion
burden~\cite{filippi_quantitative_1996,wang_survey_1997}. The number of
lesions and the appearance of new lesions on follow-up scans are central to
diagnosis and to the monitoring of disease activity, so a small lesion that
is never detected is a clinically relevant error.

Deep learning methods, particularly convolutional neural networks (CNNs),
have substantially advanced medical image segmentation. Architectures
such as U-Net and its 3D extensions have become standard for volumetric
analysis~\cite{navab_u-net_2015}, while automated frameworks such as
nnU-Net have demonstrated strong performance across a wide range of
biomedical datasets through systematic pipeline
adaptation~\cite{isensee_nnu-net_2021}. More recently, transformer-based
and hybrid CNN--transformer models have further expanded the design space
for medical segmentation~\cite{hatamizadeh_unetr_2021,
hatamizadeh_swin_2022}. However, despite these advances, segmentation of
small pathological structures remains challenging even for
state-of-the-art models: in the nnU-Net-based FLAMeS, for example, most
missed lesions were smaller than
10~mm$^3$~\cite{dereskewicz_flames_2025}.

A primary difficulty lies in the training objective. In voxel-wise
optimization, background and large lesions dominate the learning signal,
so small lesions have little influence on the learned representation. A
model trained with a voxel-wise overlap objective can therefore reach a
high Dice similarity coefficient (DSC) while leaving many small lesions
undetected. Existing losses for class imbalance reweight voxels by class,
difficulty, or position, and instance-level losses balance lesions of
different sizes, but each lesion is still assessed through overlap.
No existing objective combines size-balanced overlap with an explicit,
overlap-independent detection term for each lesion
(\cref{sec:related}).

In this work, rather than adding architectural complexity to the strong
nnU-Net baseline, we redesign the objective to give sparse lesions greater
influence during training. The proposed loss, CATMIL, adds two terms to
the standard nnU-Net loss: a \emph{Component-Adaptive Tversky (CAT)} term,
which gives small and large lesions comparable weight, and a lesion-level
\emph{Multiple Instance Learning (MIL)} term, which penalizes any lesion in
which no voxel is detected. Its practical aim is to reduce false negatives
without changing the segmentation architecture; performance is therefore
assessed with both voxel- and lesion-level metrics, on the MSLesSeg dataset
and on a second dataset, 3D-MR-MS, trained with the same loss weights.

The main contributions of this study are threefold:
\begin{itemize}
\item We propose CATMIL, a loss objective that targets the omission of
small lesions without any change to the network architecture. Because it
acts only on the training objective, it can be applied to existing
segmentation pipelines. We show analytically that the CAT term gives each
lesion a total weight and gradient that are nearly independent of its
size, and that the MIL term gives every lesion the same detection signal.
\item We show that CATMIL reduces the omission of small lesions. The
improvement is concentrated in the lesion sizes at which the standard
losses fail, it is a net gain in detected lesions rather than a trade
between lesion sizes, and it holds across test patients, across trained
models, on a second dataset, and for lesions above the size used in
clinical reading. On MSLesSeg, CATMIL also preserves the voxel-level
overlap and boundary accuracy of the standard losses. An ablation shows
that the lesion-level detection term accounts for the gain.
\item We identify why standard losses miss small lesions. The baselines do
not miss them by a narrow margin: they produce no response to most of
them, so no choice of decision threshold can recover these lesions. We
also characterize the cost of the increased sensitivity, a larger number
of small false-positive components, and show that a simple
component-size filter controls it while keeping the sensitivity gain: at
the lesion-wise precision of the standard nnU-Net loss, CATMIL detects more
small lesions.
\end{itemize}
